% ==========================================
% ФАЙЛ: p11.tex
% БЛОК БИЛЕТОВ: Степенные ряды, Евклидовы пространства и ряды Фурье
% ==========================================

% --- Вопрос 51 ---
\examquestion{51}{Коэффициенты Тейлора и единственность разложения}
{Сформулируйте и докажите теорему о связи коэффициентов сходящегося степенного ряда с производными его суммы. Проведите вывод формулы $a_n = \frac{f^{(n)}(x_0)}{n!}$ путем последовательного $n$-кратного почленного дифференцирования степенного ряда внутри его интервала сходимости и последующего зануления переменной. Из этой теоремы докажите единственность представления функции степенным рядом.}

\paragraph{Суть на пальцах}
Тейлор говорит: «Если функция совпадает с многочленом во всех своих производных в одной конкретной точке, то это одна и та же функция (в пределах радиуса сходимости)». 
Как найти $n$-й коэффициент? Мы просто $n$ раз «снимаем слои» с ряда с помощью производной. Каждая производная сбрасывает степень вниз, и к $n$-му шагу нужный нам икс остается без степени. Затем мы подставляем в ряд $x = x_0$ — это уничтожает все «хвосты» с высшими степенями. Откуда берется факториал? От правила дифференцирования: $(x^3)' = 3x^2$, $(3x^2)' = 3 \cdot 2 \cdot x = 3!x$.


\paragraph{Теорема о коэффициентах Тейлора}
Если на интервале $(x_0 - R, x_0 + R)$ функция $f(x)$ является суммой степенного ряда $f(x) = \sum_{n=0}^\infty a_n (x-x_0)^n$, то коэффициенты этого ряда однозначно определяются формулой:
$$ a_n = \frac{f^{(n)}(x_0)}{n!} $$

\begin{proof}[Вывод формулы]
1. Распишем исходный степенной ряд:
$$ f(x) = a_0 + a_1(x-x_0) + a_2(x-x_0)^2 + a_3(x-x_0)^3 + \dots $$
2. Подставим в него $x = x_0$. Все скобки $(x_0-x_0)$ зануляются, остается только первый член: $f(x_0) = a_0$, откуда $a_0 = \frac{f^{(0)}(x_0)}{0!}$.
3. Внутри интервала сходимости степенной ряд можно дифференцировать почленно бесконечное число раз. Возьмем первую производную:
$$ f'(x) = a_1 + 2a_2(x-x_0) + 3a_3(x-x_0)^2 + \dots $$
Подставим $x = x_0$: зануляется все, кроме первого члена. Получаем $f'(x_0) = a_1$.
4. Продифференцируем ряд $k$ раз. Первые $k-1$ слагаемых превратятся в нуль, а $k$-е слагаемое $a_k(x-x_0)^k$ потеряет иксы и накопит множители $k \cdot (k-1) \dots 1 = k!$. Остальные слагаемые будут содержать множитель $(x-x_0)$ в какой-то степени:
$$ f^{(k)}(x) = k! \cdot a_k + \frac{(k+1)!}{1!} a_{k+1}(x-x_0) + \frac{(k+2)!}{2!} a_{k+2}(x-x_0)^2 + \dots $$
5. Подставляя $x = x_0$, мы снова обнуляем весь бесконечный хвост:
$$ f^{(k)}(x_0) = k! \cdot a_k \implies a_k = \frac{f^{(k)}(x_0)}{k!} $$
\end{proof}

\begin{proof}[Доказательство единственности]
Предположим, что функция $f(x)$ раскладывается в два разных степенных ряда: $f(x) = \sum a_n(x-x_0)^n$ и $f(x) = \sum b_n(x-x_0)^n$. По доказанной выше теореме, коэффициенты первого ряда обязаны вычисляться по формуле $a_n = \frac{f^{(n)}(x_0)}{n!}$. Аналогично, коэффициенты второго ряда $b_n = \frac{f^{(n)}(x_0)}{n!}$. Так как правые части равны (производные функции в точке однозначны), то $a_n = b_n$ для всех $n$. Ряд единственен.
\end{proof}
\begin{proof}[]
Думайте,
\end{proof}
\begin{proof}[]
думайте.
\end{proof}



% --- Вопрос 51 ---
\begin{tcolorbox}[colback=gray!5, colframe=black, boxrule=0.5pt, arc=0pt]
\textbf{Резюмируем Вопрос 51 (Коэффициенты Тейлора и единственность):} \\
\textbf{Тезис:} Если функция представляется степенным рядом $\sum a_n(x-x_0)^n$, её коэффициенты жестко заданы как $a_n = \frac{f^{(n)}(x_0)}{n!}$, что гарантирует единственность разложения. \\
\textbf{Ход:} Дифференцируем ряд $k$ раз: первые $k-1$ членов исчезают, у $k$-го члена сбрасывается степень и накапливается факториал ($k! \cdot a_k$), а старшие члены сохраняют скобку $(x-x_0)$. \\
\textbf{Трюк:} Подстановка $x = x_0$ мгновенно обнуляет весь бесконечный старший «хвост» ряда, оставляя точное равенство $f^{(k)}(x_0) = k! \cdot a_k$. \\
\textbf{Финал:} Так как производная в точке однозначна, любые два совпадающих ряда обязаны иметь одинаковые коэффициенты ($a_n = b_n$).
\end{tcolorbox}





% --- Вопрос 52 ---
\examquestion{52}{Критерий аналитичности и ограничение роста производных}
{Сформулируйте и докажите теорему о достаточном условии аналитичности функции на промежутке: если в некоторой окрестности точки $x_0$ все производные функции ограничены единой геометрической шкалой $|f^{(n)}(x)| \le C \cdot M^n$ (где $C, M > 0$), то ряд Тейлора этой функции гарантированно сходится к ней самой. В ходе доказательства примените форму Лагранжа для остаточного члена и обоснуйте его стремление к нулю.}

\paragraph{Суть на пальцах}
Быть бесконечно дифференцируемой функцией — недостаточно, чтобы ряд Тейлора работал. Есть «больные» функции (как $e^{-1/x^2}$), у которых все производные в нуле равны нулю, ряд состоит из одних нулей, а сама функция не ноль. Чтобы ряд Тейлора точно сошелся к \textit{самой функции}, её остаток («хвост» из отброшенных бесконечных слагаемых) должен таять до нуля. Если производные не растут слишком дико (не быстрее, чем геометрическая прогрессия $M^n$), то мощный факториал $n!$ в знаменателе хвоста Тейлора их с легкостью раздавит в ноль.

\paragraph{Формулировка (Критерий аналитичности)}
Пусть функция $f(x)$ бесконечно дифференцируема в окрестности $U_{\delta}(x_0)$. Если существуют такие константы $C > 0$ и $M > 0$, что для всех $x \in U_{\delta}(x_0)$ и для всех $n \in \mathbb{N}$ выполняется оценка:
$$ |f^{(n)}(x)| \le C \cdot M^n $$
то функция $f(x)$ является аналитической в этой окрестности, то есть её ряд Тейлора сходится к ней самой: $f(x) = \sum_{n=0}^\infty \frac{f^{(n)}(x_0)}{n!} (x-x_0)^n$.

\begin{proof}
1. Запишем формулу Тейлора для функции $f(x)$ с остаточным членом $R_n(x)$ в форме Лагранжа:
$$ f(x) = S_n(x) + R_n(x) $$
где $S_n(x)$ — $n$-я частичная сумма ряда Тейлора, а остаточный член имеет вид:
$$ R_n(x) = \frac{f^{(n+1)}(\xi)}{(n+1)!} (x-x_0)^{n+1} $$
где $\xi$ — некоторая точка между $x_0$ и $x$.
2. Наша цель — доказать, что при $n \to \infty$ остаток $R_n(x) \to 0$. Перейдем к модулям и применим условие теоремы к производной порядка $n+1$:
$$ |R_n(x)| = \frac{|f^{(n+1)}(\xi)|}{(n+1)!} |x-x_0|^{n+1} \le \frac{C \cdot M^{n+1}}{(n+1)!} |x-x_0|^{n+1} $$
3. Объединим степени в одну скобку:
$$ |R_n(x)| \le C \cdot \frac{(M |x-x_0|)^{n+1}}{(n+1)!} $$
4. Обозначим фиксированное неотрицательное число $A = M |x-x_0|$. Известный предел матанализа (доказываемый через признак Д'Аламбера) утверждает, что факториал растет быстрее любой показательной функции, поэтому для любого фиксированного $A$:
$$ \lim_{n \to \infty} \frac{A^{n+1}}{(n+1)!} = 0 $$
5. Следовательно, $\lim_{n \to \infty} R_n(x) = 0$.
6. Перейдем к пределу в исходной формуле Тейлора:
$$ \lim_{n \to \infty} S_n(x) = f(x) - \lim_{n \to \infty} R_n(x) = f(x) - 0 = f(x) $$
Частичные суммы ряда Тейлора сходятся к $f(x)$. Теорема доказана.
\end{proof}
Итог: Мы не суммируем хвост. Мы доказываем, что размер всей этой кучи (хвоста) не может превысить значение $R_n(x)$. А так как $R_n(x) \to 0$ при $n \to \infty$, то и вся куча (бесконечный хвост) обязана схлопнуться в ноль.
% --- Вопрос 52 ---
\begin{tcolorbox}[colback=gray!5, colframe=black, boxrule=0.5pt, arc=0pt]
\textbf{Резюмируем Вопрос 52 (Критерий аналитичности):} \\
\textbf{Тезис:} Ряд Тейлора гарантированно сходится к самой функции, если её производные ограничены единой геометрической шкалой: $|f^{(n)}(x)| \le C \cdot M^n$. \\
\textbf{Ход:} Записываем остаток в форме Лагранжа: $R_n(x) = \frac{f^{(n+1)}(\xi)}{(n+1)!} (x-x_0)^{n+1}$ и навешиваем модули. \\
\textbf{Трюк:} Оцениваем числитель через шкалу: $|R_n(x)| \le C \cdot \frac{(M|x-x_0|)^{n+1}}{(n+1)!}$. Фиксируем числитель как $A^{n+1}$. \\
\textbf{Финал:} Так как факториал растет быстрее экспоненты ($\lim \frac{A^{n+1}}{(n+1)!} = 0$), остаточный член $R_n(x) \to 0$. Ряд Тейлора «схлопывается» точно в функцию.
\end{tcolorbox}



% --- Вопрос 53 ---
\examquestion{53}{Евклидовы пространства функций и неравенство Коши-Буняковского}
{Сформулируйте понятие предгильбертова пространства функций на базе интеграла Римана. Проведите доказательство неравенства Коши-Буняковского, рассматривая дискриминант неотрицательного квадратичного трехчлена $||\lambda f+g||^2 \ge 0$. Обоснуйте структуру фактор-пространства для обеспечения корректности метрики.}

\paragraph{Суть на пальцах}
Мы берем школьную геометрию 3D-пространства (длины и углы) и натягиваем её на мир функций. Вместо скалярного умножения векторов $(a_1b_1 + a_2b_2)$ мы используем интеграл произведения $\int f(x)g(x)dx$. Коши-Буняковский просто говорит, что косинус угла между любыми функциями не больше единицы. 
Но есть математический «баг»: функция может быть всюду нулем, но в одной точке прыгнуть в бесконечность. Площадь (интеграл) под ней будет ноль, и длина будет ноль, хотя функция не пустая. Чтобы починить это, математики договорились: если функции отличаются друг от друга только в наборе изолированных точек (на множестве нулевой площади), мы считаем их одной и той же функцией.

\paragraph{Определение предгильбертова пространства функций}
Линейное пространство вещественных функций называется предгильбертовым (или евклидовым), если на нем задано скалярное произведение со свойствами билинейности, симметричности и положительной определенности. На базе интеграла Римана для $f, g \in \mathcal{R}[a, b]$ оно вводится как:
$$ (f, g) = \int_a^b f(x)g(x) dx $$
Норма (длина) функции определяется как $||f|| = \sqrt{(f, f)}$.

\paragraph{Структура фактор-пространства}
Аксиома нормы требует: $||f|| = 0 \iff f(x) \equiv 0$. Однако для интеграла Римана $\int_a^b f^2(x) dx = 0$ может выполняться и для функции, отличной от тождественного нуля (например, ненулевой лишь в одной точке). Для корректности метрики вводится отношение эквивалентности: $f \sim g$, если $f(x) = g(x)$ почти всюду (на множестве Лебеговой меры нуль). Пространство строится не из самих функций, а из \textbf{классов эквивалентности} таких функций. В таком фактор-пространстве аксиома $||f|| = 0 \iff f = [0]$ выполняется строго.

\paragraph{Неравенство Коши-Буняковского}
Для любых $f, g$ из евклидова пространства выполняется:
$$ |(f, g)| \le ||f|| \cdot ||g|| $$

\begin{proof}
1. Рассмотрим линейную комбинацию функций $\lambda f + g$, где $\lambda \in \mathbb{R}$ — произвольная константа.
2. По аксиомам нормы квадрат нормы любой функции всегда неотрицателен:
$$ ||\lambda f + g||^2 \ge 0 $$
3. Раскроем квадрат нормы через свойства скалярного произведения (линейность и симметрию):
$$ (\lambda f + g, \lambda f + g) = \lambda^2(f, f) + 2\lambda(f, g) + (g, g) \ge 0 $$
4. Перепишем это через нормы:
$$ ||f||^2 \lambda^2 + 2(f, g)\lambda + ||g||^2 \ge 0 $$
5. Если $||f|| = 0$, то $f \sim 0$, тогда $(f, g) = 0$, и неравенство $|0| \le 0 \cdot ||g||$ выполняется тривиально. 
6. Пусть $||f|| > 0$. Тогда перед нами квадратичный трехчлен относительно $\lambda$ вида $A\lambda^2 + 2B\lambda + C \ge 0$, где $A = ||f||^2 > 0$.
7. Так как парабола ветвями вверх всегда неотрицательна (находится выше или касается оси $\lambda$), её дискриминант (или четверть дискриминанта) обязан быть неположительным:
$$ \frac{D}{4} = B^2 - AC \le 0 $$
8. Подставляем коэффициенты:
$$ (f, g)^2 - ||f||^2 \cdot ||g||^2 \le 0 \implies (f, g)^2 \le ||f||^2 ||g||^2 $$
9. Извлекая корень из обеих частей, получаем требуемое: $|(f, g)| \le ||f|| \cdot ||g||$.
\end{proof}

% --- Вопрос 53 ---
\begin{tcolorbox}[colback=gray!5, colframe=black, boxrule=0.5pt, arc=0pt]
\textbf{Резюмируем Вопрос 53 (Евклидовы пространства и Коши-Буняковский):} \\
\textbf{Тезис:} Скалярное произведение функций задается как $(f,g) = \int_a^b f(x)g(x)dx$. Неравенство $|(f,g)| \le \|f\| \cdot \|g\|$ гарантирует корректность геометрии. \\
\textbf{Ход:} Составляем заведомо неотрицательный квадрат нормы линейной комбинации: $\|\lambda f + g\|^2 = \lambda^2 \|f\|^2 + 2\lambda (f,g) + \|g\|^2 \ge 0$. \\
\textbf{Трюк:} Полученный квадратный трехчлен относительно $\lambda$ не может иметь двух вещественных корней $\implies$ его дискриминант $D/4 = (f,g)^2 - \|f\|^2\|g\|^2 \le 0$. \\
\textbf{Финал:} Извлечение корня дает неравенство Коши-Буняковского. Для починки аксиомы нормы ($\|f\|=0 \iff f=0$) функции объединяют в классы эквивалентности (фактор-пространство по функциям, равным 0 почти всюду).
\end{tcolorbox}


% --- Вопрос 54 ---
\examquestion{54}{Экстремальное свойство коэффициентов Фурье}
{Сформулируйте и докажите теорему об экстремальном свойстве коэффициентов Фурье (минимизация среднеквадратичного отклонения линейной комбинации от функции). Из доказательства выведите неравенство Бесселя.}

\paragraph{Введение в курс дела: База про пространство и коэффициенты}
Прежде чем доказывать, зафиксируем правила игры. Мы находимся в евклидовом пространстве функций, которое работает по законам обычной геометрии:
\begin{itemize}
    \item \textbf{Скалярное произведение и норма:} Норма (длина вектора/функции) в квадрате — это скалярное произведение элемента на самого себя: $||x||^2 = (x, x)$. Скалярное произведение обладает свойством линейности, то есть скобки можно раскрывать как в обычной алгебре: $(x-y, x-y) = (x, x) - 2(x, y) + (y, y)$.
    \item \textbf{Ортогональность:} Базис $\{e_k\}$ ортонормирован. Это значит, что длина каждого базисного вектора равна 1 ($(e_k, e_k) = 1$), а любые два разных вектора перпендикулярны, их скалярное произведение равно нулю ($(e_i, e_j) = 0$). Благодаря этому при возведении больших сумм в квадрат все "перекрестные" слагаемые просто исчезают.
    \item \textbf{Кто такие $c_k$ и $a_k$?} $c_k$ — это абсолютно \textit{любые} числа-коэффициенты. Мы берем их наугад, пытаясь собрать из базисных функций $e_k$ некую комбинацию (многочлен), которая была бы максимально похожа на нашу функцию $f$. А вот $a_k$ — это строгие \textbf{коэффициенты Фурье}, которые вычисляются как точная проекция функции на конкретную базисную ось: $a_k = (f, e_k)$. 
\end{itemize}
Вся суть теоремы сводится к тому, чтобы доказать: наша попытка подогнать "левые" коэффициенты $c_k$ увенчается максимальным успехом (ошибка приближения будет минимальной) только тогда, когда мы перестанем гадать и просто приравняем $c_k$ к истинным коэффициентам Фурье $a_k$.

\paragraph{Формулировка теоремы}
Пусть $\{e_k\}$ — ортонормированная система в евклидовом пространстве. Среди всех возможных линейных комбинаций вида $P_n = \sum_{k=1}^n c_k e_k$, наименьшее среднеквадратичное отклонение от функции $f$ достигается тогда и только тогда, когда коэффициенты $c_k$ равны коэффициентам Фурье этой функции: $c_k = a_k = (f, e_k)$.

\begin{proof}
1. Запишем квадрат среднеквадратичного отклонения (ошибки $\Delta^2$):
$$ \Delta^2 = \left|\left|f - \sum_{k=1}^n c_k e_k\right|\right|^2 $$
2. Раскроем квадрат через скалярное произведение (по правилу $(A-B, A-B)$):
$$ \Delta^2 = \left(f - \sum_{k=1}^n c_k e_k, f - \sum_{k=1}^n c_k e_k\right) = (f, f) - 2\sum_{k=1}^n c_k(f, e_k) + \left(\sum_{k=1}^n c_k e_k, \sum_{k=1}^n c_k e_k\right) $$
3. Воспользуемся свойством ортонормированности: $(e_i, e_j) = 0$ при $i \ne j$, и $(e_k, e_k) = 1$. Тогда громоздкое скалярное произведение двух сумм превращается в простую сумму квадратов:
$$ \Delta^2 = ||f||^2 - 2\sum_{k=1}^n c_k (f, e_k) + \sum_{k=1}^n c_k^2 $$
4. По определению, коэффициент Фурье $a_k = (f, e_k)$. Подставим это обозначение во второе слагаемое:
$$ \Delta^2 = ||f||^2 - 2\sum_{k=1}^n c_k a_k + \sum_{k=1}^n c_k^2 $$
5. Искусственно добавим и вычтем сумму квадратов коэффициентов Фурье $\sum_{k=1}^n a_k^2$, чтобы свернуть выражение в полный квадрат:
$$ \Delta^2 = ||f||^2 - \sum_{k=1}^n a_k^2 + \sum_{k=1}^n (c_k^2 - 2c_k a_k + a_k^2) = ||f||^2 - \sum_{k=1}^n a_k^2 + \sum_{k=1}^n (c_k - a_k)^2 $$
6. Первые два слагаемых в правой части зависят только от самой функции $f$ и её коэффициентов $a_k$ (они жестко фиксированы функцией). Изменить нашим выбором можно только третье слагаемое (оно всегда $\ge 0$, так как это сумма квадратов). 
Очевидно, что общая ошибка $\Delta^2$ будет минимальной, когда сумма квадратов $\sum (c_k - a_k)^2$ обратится в строгий ноль. Это происходит тогда и только тогда, когда для всех $k$ выполняется равенство: $c_k = a_k$.
\end{proof}

\paragraph{Вывод неравенства Бесселя}
Возьмем доказанную формулу для ошибки $\Delta^2$ в случае идеального выбора коэффициентов ($c_k = a_k$). Третье слагаемое зануляется:
$$ \Delta_{min}^2 = ||f||^2 - \sum_{k=1}^n a_k^2 $$
Так как квадрат любой метрики неотрицателен по определению, $\Delta_{min}^2 \ge 0$. Следовательно:
$$ ||f||^2 - \sum_{k=1}^n a_k^2 \ge 0 \implies \sum_{k=1}^n a_k^2 \le ||f||^2 $$
Это неравенство Бесселя. При переходе к пределу при $n \to \infty$ оно доказывает сходимость числового ряда из квадратов коэффициентов Фурье: $\sum_{k=1}^\infty a_k^2 \le ||f||^2$.


% --- Вопрос 54 ---
\begin{tcolorbox}[colback=gray!5, colframe=black, boxrule=0.5pt, arc=0pt]
\textbf{Резюмируем Вопрос 54 (Экстремальное свойство и Бессель):} \\
\textbf{Тезис:} Наименьшая среднеквадратичная ошибка приближения функции комбинацией $\sum c_k e_k$ достигается только при равенстве $c_k$ коэффициентам Фурье $a_k = (f,e_k)$. \\
\textbf{Ход:} Раскрываем квадрат ошибки $\|f - \sum c_k e_k\|^2$ через скалярное произведение, используя ортонормированность базиса ($(e_i,e_j)=\delta_{ij}$). \\
\textbf{Трюк:} Искусственно выделяем полный квадрат относительно коэффициентов: $\Delta^2 = \|f\|^2 - \sum a_k^2 + \sum (c_k - a_k)^2$. \\
\textbf{Финал:} Минимум $\Delta^2$ достигается при занулении последней суммы ($c_k = a_k$). Так как $\Delta_{min}^2 \ge 0$, автоматически вылетает неравенство Бесселя: $\sum a_k^2 \le \|f\|^2$.
\end{tcolorbox}





% --- Вопрос 55 ---
\examquestion{55}{Теорема Стеклова}
{Сформулируйте понятия полноты и замкнутости системы функций и докажите их эквивалентность в гильбертовом пространстве.}

\paragraph{Суть на пальцах}
Замкнутая система — это такая, которой «хватает мощности», чтобы описать любой сигнал без остатка. Вся энергия функции без потерь рассыпается на коэффициенты (выполняется точное равенство Парсеваля, линейная оболочка плотна). Полная система — это такая, где нет «слепых зон»: если какой-то сигнал перпендикулярен (ортогонален) абсолютно всем базисным векторам, то это может быть только нулевой сигнал. Стеклов доказал, что в гильбертовом пространстве (где нет дырок) эти два понятия — одно и то же.

\paragraph{Определения}
Пусть $\{e_k\}$ — ортонормированная система в гильбертовом пространстве $H$.
\begin{itemize}
    \item \textbf{Замкнутость:} Линейная оболочка $\{e_k\}$ плотна в $H$. То есть для любой $f \in H$ и любого $\varepsilon > 0$ существуют коэффициенты $c_1, \dots, c_N$ такие, что $||f - \sum_{k=1}^N c_k e_k|| < \varepsilon$.
    \item \textbf{Полнота:} Не существует ненулевого элемента $f \in H$, ортогонального всем $e_k$. То есть: $(\forall k: (f, e_k) = 0) \implies f = 0$.
\end{itemize}

\paragraph{Теорема Стеклова}
В гильбертовом пространстве $H$ ортонормированная система $\{e_k\}$ является полной тогда и только тогда, когда она является замкнутой.

\begin{proof}
\textbf{1. Замкнутость $\implies$ Полнота}
1. Берем произвольную функцию $f \in H$, ортогональную всем базисным векторам: $(f, e_k) = 0$ для всех $k$.
2. По условию замкнутости (плотности), для любого $\varepsilon > 0$ мы можем приблизить $f$ линейной комбинацией $S = \sum_{k=1}^N c_k e_k$ так, что $||f - S|| < \varepsilon$.
3. Используем ортогональность. Так как $f$ ортогональна каждому $e_k$, она ортогональна и их линейной комбинации $S$, то есть $(f, S) = 0$.
Распишем квадрат нормы $f$:
$$ ||f||^2 = (f, f) = (f, f - S) + (f, S) = (f, f - S) $$
4. Применяем неравенство Коши–Буняковского к скалярному произведению $(f, f - S)$:
$$ ||f||^2 = |(f, f - S)| \le ||f|| \cdot ||f - S|| $$
5. Если $||f|| = 0$, то $f = 0$, и полнота доказана.
Предположим противное: $||f|| > 0$. Тогда разделим обе части на $||f||$:
$$ ||f|| \le ||f - S|| $$
6. Но по условию замкнутости мы можем выбрать $S$ так, чтобы разность $||f - S||$ была \textbf{сколь угодно малой}. Выберем её строго меньше, чем $||f||/2$. Тогда:
$$ ||f|| \le ||f - S|| < \frac{||f||}{2} $$
Получили противоречие ($1 \le 1/2$). Значит, единственный выход: $||f|| = 0 \implies f = 0$. Полнота доказана.

\textbf{2. Полнота $\implies$ Замкнутость}
1. Пусть система полна. Возьмем любую $f \in H$ и её коэффициенты Фурье $a_k = (f, e_k)$. 
2. По неравенству Бесселя ряд $\sum a_k^2$ сходится. В силу полноты самого гильбертова пространства (теорема Рисса-Фишера), ряд $\sum a_k e_k$ сходится к некоторому элементу $\tilde{f} \in H$.
3. Рассмотрим ошибку: $g = f - \tilde{f}$. Вычислим её скалярное произведение с любым базисным вектором $e_m$:
$$ (g, e_m) = (f - \tilde{f}, e_m) = (f, e_m) - (\tilde{f}, e_m) = a_m - a_m = 0 $$
4. Вектор $g$ ортогонален всем $e_k$. Но система полна, значит, нулевым может быть только нулевой вектор. Отсюда $g = 0 \implies f = \tilde{f}$.
5. Ряд Фурье сходится к $f$, линейная оболочка плотна. Замкнутость доказана.
\end{proof}

\paragraph{Короткая фраза для защиты (Замкнутость $\implies$ Полнота)}
\begin{quote}
«Если система замкнута, любую функцию можно сколь угодно точно приблизить комбинациями $e_k$. Если взять функцию, ортогональную всему базису, то по Коши-Буняковскому её норма не превзойдет нормы разности с этим приближением. Делая приближение бесконечно точным, мы сжимаем норму функции до нуля. Значит, функция нулевая — это и есть полнота.»
\end{quote}

% --- Вопрос 55 ---
\begin{tcolorbox}[colback=gray!5, colframe=black, boxrule=0.5pt, arc=0pt]
\textbf{Резюмируем Вопрос 55 (Теорема Стеклова):} \\
\textbf{Тезис:} В гильбертовом пространстве понятия полноты (нет векторов, ортогональных всему базису, кроме нуля) и замкнутости (оболочка плотна) эквивалентны. \\
\textbf{Ход (Замкнутость $\implies$ Полнота):} Пусть $(f, e_k) = 0 \,\, \forall k$. По замкнутости приближаем $f$ линейной комбинацией $S$ с точностью $\varepsilon$. \\
\textbf{Трюк:} Так как $(f, S) = 0$, расписываем $\|f\|^2 = (f, f-S) \le \|f\| \cdot \|f-S\|$ по Коши-Буняковскому. Деля на $\|f\|$, получаем $\|f\| \le \|f-S\| < \varepsilon$. \\
\textbf{Финал:} Устремляя $\varepsilon \to 0$, зажимаем норму функции в ноль ($\|f\| = 0 \implies f = 0$). Обратная сторона доказывается через сходимость ряда Фурье по Риссу-Фишеру к элементу $\tilde{f}$ и обнуление ошибки $f - \tilde{f}$.
\end{tcolorbox}